Der Alkohol sei Ethanol mit $L_\text{v} = \LvEO$ und $c = \cEtO$ (FoTaBe S.~191). Verdunstet der Anteil $f$ der Masse $m$ in kurzer Zeit, so stammt die Verdampfungswärme aus der restlichen Flüssigkeit (Wärme von Schale und Luft vernachlässigt):
\begin{align}
 f\,m\,L_\text{v} &= \left(1-f\right)m\,c\,\Delta\vartheta
\end{align}
Die Masse kürzt sich weg, und die Temperatur sinkt um
\begin{align}
 \Delta\vartheta &= \dTF \\
   &= \frac{\fvO\cdot\LvE}{\left(1-\fvO\right)\cdot\cEt} \\
   &= \dT \approx \result{\dTP{0}{2}}
\end{align}
Die Tabelle gibt die Verdampfungswärme beim Siedepunkt; bei Zimmertemperatur ist sie etwas grösser, die Abkühlung also eher noch etwas stärker.
